\subsubsection{Sweep line template}

\paragraph{Idea intuitiva.}
Patron de diseno para problemas geometricos donde se barre una linea sobre el plano y se mantiene una estructura de datos sobre el otro eje. Eventos (add, query, remove) se ordenan por la coordenada de barrido y se procesan secuencialmente.

\paragraph{Propiedades clave (invariantes).}
\begin{itemize}\setlength\itemsep{0pt}
	\item Eventos: tuplas \texttt{(x, type, id)}.
	\item La estructura auxiliar (set, multiset, segtree) se mantiene sobre el otro eje.
	\item Procesar eventos en orden de $x$ ascendente.
\end{itemize}

\paragraph{Complejidad.}
\begin{itemize}\setlength\itemsep{0pt}
	\item Depende del problema; tipico $O((N + E) \log N)$.
\end{itemize}

\paragraph{Donde tocar para modificar.}
\begin{itemize}\setlength\itemsep{0pt}
	\item \textbf{Aniadir/quitar segmentos}: aniadir eventos add/remove al barrido.
	\item \textbf{Aniadir queries}: aniadir eventos query.
	\item \textbf{Cambiar el eje de barrido}: si barres en $y$ en vez de $x$, aniadir comparador.
\end{itemize}

\paragraph{Trampas y errores frecuentes.}
\begin{itemize}\setlength\itemsep{0pt}
	\item El snippet es solo el \textbf{esqueleto}; los handlers \texttt{add/query/remove} se llenan por problema.
	\item Cuidado con eventos duplicados.
\end{itemize}

\paragraph{Codigo.}
Ver \texttt{cpp.tex}, seccion \textit{Sweep line template}.

\vspace{0.5em}
